\documentclass[11pt]{article}

\usepackage[margin=1in]{geometry}
\usepackage{amsmath, amssymb}
\usepackage{xcolor}
\usepackage{listings}
\usepackage{hyperref}
\usepackage{array}
\usepackage{booktabs}

\hypersetup{
    colorlinks=true,
    linkcolor=blue!60!black,
    urlcolor=blue!60!black,
    citecolor=blue!60!black
}

\lstdefinestyle{bugpython}{
  language=Python,
  basicstyle=\ttfamily\small,
  keywordstyle=\color{blue!60!black},
  stringstyle=\color{green!40!black},
  commentstyle=\color{gray!70!black},
  showstringspaces=false,
  breaklines=true,
  frame=single,
  rulecolor=\color{gray!30},
  columns=fullflexible,
  keepspaces=true
}

\title{SymPy Bug Report}
\author{}
\date{}

\begin{document}

\maketitle

\section{Bug 3: solveset returns \texttt{nan} as a zero of secant}

\subsection{Minimal Reproducer}

\medskip

\begin{lstlisting}[style=bugpython]
from sympy import S, cos, sec, solveset, symbols

x = symbols("x")

solution = solveset(sec(x), x, domain=S.Complexes)
print("solveset(sec(x), x, domain=S.Complexes) =", solution)
print("sec(nan) =", sec(S.NaN))
print("solveset(1/cos(x), x, domain=S.Complexes) =",
      solveset(1 / cos(x), x, domain=S.Complexes))
\end{lstlisting}

\subsection{Observed Behavior}

\medskip

\begin{lstlisting}[style=bugpython]
SymPy version: 1.14.0
solveset(sec(x), x, domain=S.Complexes) = {nan}
sec(nan) = nan
solveset(1/cos(x), x, domain=S.Complexes) = EmptySet
\end{lstlisting}

SymPy returns a singleton set containing \(\mathtt{nan}\) as a zero of
\(\sec(x)\). Direct evaluation of the reported candidate gives
\(\sec(\mathtt{nan})=\mathtt{nan}\), not \(0\), and the mathematically
equivalent reciprocal spelling \(1/\cos(x)\) returns \(\varnothing\).

\subsection{Expected Behavior}

The expected result is
\[
    \operatorname{solveset}(\sec(x), x, \mathbb{C}) = \varnothing.
\]
The value \(\mathtt{nan}\) is not a complex number and should never be returned
as a solution of the original equation.

\subsection{Mathematical Explanation}

For every finite complex number \(z\), \(\cos(z)\) is finite. Therefore
\[
    \sec(z) = \frac{1}{\cos(z)}
\]
can never equal \(0\). If \(\cos(z)\neq 0\), then \(\sec(z)\) is a nonzero
finite reciprocal. If \(\cos(z)=0\), then \(\sec(z)\) has a pole and is
undefined or infinite, not zero. Thus there is no finite complex \(z\) satisfying
\[
    \sec(z)=0.
\]
The returned value is not a harmless alternative representation or a branch
choice: \(\mathtt{nan}\) is not an element of the complex solution domain, and
substitution does not satisfy the equation.

\subsection{Source-Level Diagnosis}

The diagnosis identifies the relevant source file as
\nolinkurl{sympy/solvers/solveset.py}. The public call
\texttt{solveset(sec(x), x, domain=S.Complexes)} reaches \texttt{\_solveset};
the fallback inversion path then calls
\texttt{invert\_complex(f, 0, symbol, domain)} around lines 2321--2324.
\texttt{invert\_complex} delegates to \texttt{\_invert\_complex}, which
recognizes \(\sec(x)\) as a trigonometric function around lines 577--578 and
then calls \texttt{\_invert\_trig\_hyp\_complex}.

The problematic rule is in \texttt{\_invert\_trig\_hyp\_complex} around lines
462--466: the inversion helper treats \(\sec\) like an ordinary invertible
trigonometric function, chooses \texttt{asec} for the inverse, and constructs
families of the form \(2n\pi+\operatorname{asec}(a)\) and
\(2n\pi-\operatorname{asec}(a)\). For the equation \(\sec(x)=0\), this applies
\(\operatorname{asec}(0)\). In SymPy 1.14.0, \texttt{asec(0)} evaluates to
\texttt{zoo}; arithmetic such as \(2n\pi\pm\texttt{zoo}\) simplifies during set
construction to \(\{\mathtt{nan}\}\). The inversion code does not first enforce
the reciprocal-trig range condition that \(\sec\) and \(\csc\) cannot take the
value \(0\), and it does not later filter non-finite or \(\mathtt{nan}\)
candidates from the returned finite set.

\begin{lstlisting}[style=bugpython]
# Simplified from the diagnosed inversion step:
# sec(x) = 0 is inverted as
x = 2*n*pi + asec(0)
x = 2*n*pi - asec(0)
\end{lstlisting}

This source-level behavior turns an impossible equation into invalid inverse
images. A plausible fix direction supported by the diagnosis is to add range and
singularity guards in the trigonometric inversion helper: for reciprocal
trigonometric functions such as \(\sec\) and \(\csc\), exclude \(0\) from the
target values before applying \texttt{asec} or \texttt{acsc}; if no target values
remain, return \texttt{S.EmptySet}. A defensive cleanup could also discard
\(\mathtt{nan}\), \(\infty\), \(-\infty\), and complex infinity from finite
solution sets.

\subsection{Additional Failing Instances}

The independent verification checks the representative case and five additional
instances of the same form,
\[
    a\,\sec(bx+c),
\]
where \(a\neq 0\). Each expression has no complex zeros for the same reason:
away from poles it is a nonzero scalar times a reciprocal of a finite cosine,
and at poles it is undefined or infinite rather than zero. SymPy nevertheless
returns \(\{\mathtt{nan}\}\) for each listed case, while the equivalent
reciprocal spelling \(a/\cos(bx+c)\) returns \(\varnothing\).

\begin{center}
\small
\renewcommand{\arraystretch}{1.3}
\begin{tabular}{@{}>{\raggedright\arraybackslash}p{0.44\linewidth}>{\raggedright\arraybackslash}p{0.23\linewidth}>{\raggedright\arraybackslash}p{0.23\linewidth}@{}}
\toprule
\textbf{Input / Expression} & \textbf{SymPy behavior} & \textbf{Expected behavior} \\
\midrule
\(\sec(x)\) & Returns \(\{\mathtt{nan}\}\) & \(\varnothing\) \\
\(\sec(x+1)\) & Returns \(\{\mathtt{nan}\}\) & \(\varnothing\) \\
\(\sec(x-2)\) & Returns \(\{\mathtt{nan}\}\) & \(\varnothing\) \\
\(\sec(2x)\) & Returns \(\{\mathtt{nan}\}\) & \(\varnothing\) \\
\(\sec(3x+1)\) & Returns \(\{\mathtt{nan}\}\) & \(\varnothing\) \\
\(2\sec(x)\) & Returns \(\{\mathtt{nan}\}\) & \(\varnothing\) \\
\bottomrule
\end{tabular}
\end{center}

\subsection{Related Issues Assessment}

The bug-finding harness performed a best-effort deduplication check and
classified this as a new instance of a known failure mode. The closest related
public material is SymPy's solveset documentation, SymPy's elementary-function
documentation for reciprocal trigonometric functions, and the SymPy 1.14.0
source code showing the \(\sec\)/\(\csc\) inversion path through
\texttt{asec}/\texttt{acsc}. The artifact did not identify a public issue, pull
request, release note, Stack Overflow report, or discussion describing this
exact return value for \texttt{solveset(sec(x), x, domain=S.Complexes)}. The
general reciprocal-trig inversion failure mode may be broadly understood, but
this exact reproducible case is previously unreported and remains individually
actionable as an issue and regression test.

\subsection{Proposed SymPy Regression Test}

\medskip

\begin{lstlisting}[style=bugpython]
import pytest

from sympy import S, sec, solveset, symbols


@pytest.mark.parametrize(
    "scale, slope, offset",
    [
        (1, 1, 0),
        (1, 1, 1),
        (1, 1, -2),
        (1, 2, 0),
        (1, 3, 1),
        (2, 1, 0),
    ],
)
def test_solveset_sec_has_no_complex_zero(scale, slope, offset):
    x = symbols("x")
    assert solveset(scale * sec(slope * x + offset), x, domain=S.Complexes) == S.EmptySet
\end{lstlisting}

\end{document}
